\subsection{Компактные операторы}
% рассказ про чипсятины Lays
\begin{definition}
    $E_1, E_2$~--- ЛНП, $A \in \L(E_1, E_2)$. $A$ называется \textit{компактным}, если $\forall$ огр. $S \in E$ $A(S)$~--- предкомпакт (в $E_2$). $K(E_1, E_2)$~--- множество компактных операторов
\end{definition}


\textbf{Напоминание:}
В банаховом пространстве предкомпактность эквивалентна вполне ограниченности

\begin{definition}[эквивалентное, <<с волной>>] % будем использовать
$A$ называется \textit{компактным}, если 
$A( \overline{B}(0, 1) )$~--- предкомпакт.

\textit{Аналогия с определением ограниченного оператора}
\end{definition}

% тут Поля должна вставить картинку про эквивалентность определений

Вполне непрерывные операторы (<<улучшающая сходимость>>):
\[
x_n \weakconv x \xrightarrow{A}[] \|Ax_n - Ax\| \to 0
\]

\begin{exercise}
    $E$~--- банахово, $A \in K(E), E \ni x_n \weakconv x \in E \implies \|Ax_n - Ax\| \to 0$
    
    Подсказка:

    \[
    \begin{cases}
    y_n \weakconv y\\
    \{y_n\} \text{~--- предкомпакт}
    \end{cases}
    \Rightarrow \|y_n - y\| \to 0
    \]
\end{exercise}

\begin{lemmanote}
    Если $E$~--- рефлексивное банахово, то компактность $\sim$ вполне непрерывность.
\end{lemmanote}

% переселились в бесконечномерные пространства

\begin{example}
Ищем первообразную для функции $f(x)$:
$y' = f(x), y(0) = 0$. Это эквивалентно $y(x)=\int_{t=0}^x f(t)dt$

Резольвента~--- это <<решатель>> задачи. В данном случае <<решателем>> является оператор Вольтерра.

% Поооооля, не хватает картиночек
\end{example}

\begin{statement}[1]
Если $\dim E_1 < \infty$ или $\dim E_2 < \infty$,
$A \in \L(E_1, E_2)$, то $A$~--- компактный оператор
\end{statement}

\begin{proof}
\begin{itemize}
    \item $\dim E_1 < \infty$
        $M \subset E_1$ ограничено, а значит вполне ограничено в банаховом пространстве. $A$~--- непрерывный, а значит
        $M \subset \overline{M}, AM \subset A\overline{M}$~--- компакт (непрерывный образ предкомпакта~--- компакт). $\overline{AM} \subset A\overline{M}$.
        Замкнутое множество внутри компакта~--- компакт
    \item $\dim E_2 < \infty$. $M \subset E_1$ ограничено, $AM$ ограничено в $E_2$, а значит вполне ограничено и является предкомпактом.
\end{itemize}
\end{proof}

\begin{example}
$E = l_2(\C)$, $I$~--- тождественный оператор.

$I(\overline{B}(0, 1)) = \overline{B}(0, 1)$~--- не вполне ограничено.
\end{example}

\begin{statement}[2]
    $K(E)$~--- идеал в $\L(E)$, то есть $\forall B \in \L(E) \: \forall A \in K(E) \implies AB, BA \in K(E)$. (Идеал в обычном алгебраическом смысле).
\end{statement}

% Мы можем поиграть в эту игру. Я выиграю. Приятно знать, что ты выиграешь

\begin{proof}
очевидно  % TODO: расписать "очевидность"
\end{proof}

\begin{statement}[3]
Если $\dim E = \infty$, то $I \not \in K(E)$ ($I$~--- тождественный оператор).
\end{statement}

\begin{note}
Хотим узнать, конечномерно ли пространство $E$?
Теперь появился новый способ описания конечномерных пространств:

$\dim E < \infty \; \text{единичная сфера вполне ограничена} \;  I\text{~--- компактный оператор}$
\end{note}

\begin{proof}
    Теорема \nameref{4.1} и доказанные утверждения.
\end{proof}

\begin{corollary}
    $\dim E = \infty, A \in K(E) \implies 0 \in \sigma(A)$
\end{corollary}

\begin{proof}
    Предположим противное, то есть $\underbrace{(A - 0 I)^{-1}}_{=A^{-1}} \in \L(E)$. $AA^{-1} = I$ % TODO: ???
\end{proof}

\begin{theorem}[12.1]
    Пусть $E_1$~--- нормированное, $E_2$~--- банахово пространство, $\forall n \: A_n \in K(E_1, E_2)$ и $\|A_n - A\| \to 0$. Утверждается, что $A \in K(E_1, E_2)$.
\end{theorem}

\begin{example}
    В гильбертовом пространстве $l_2$ рассматривается диагональный оператор
    $Ax = (\lambda x_1, \lambda x_2, \ldots)$, $\{\lambda_n\}$~--- собственные числа оператора $A$. Этот оператор будет компактным тогда и только тогда, когда $\lambda_n \to 0$.
    
    Решается с помощью операторов проектирования и показывается, что последовательность таких операторов по норме стремится к 0.
\end{example}

В 1973 году доказали, что не всегда компактные операторы представляются как предел операторов конечного ранга.

% mathnet.ru - самое важное что мы должны узнать из лекции

\begin{proof}
    Подсказка самим себе: аналогия с равномерной сходимостью непрерывных функций

    <<Тут картинка>> % TODO

\begin{itemize}
    \item[1.] Нужно доказать, что $A(\overline{B})$~--- замкнутый ($\overline{B}$ - единичный шар из определения <<с волной>>). Зафиксируем эпсилон из определения предела
    \[ \forall \varepsilon > 0 \: \exists N: \: \forall n \geqslant N \: \|A_n-A\| < \varepsilon\]
    % обводим себя в кружок чтобы мы никуда не убегали
    \item[2.] Фиксируем $n \geqslant N$
    $A_n \overline{B}$~--- вполне ограничено, следовательно
    $\exists \varepsilon$-сеть $\{y_1^n, \ldots y_{m_n}^m\}$ для $A_n \overline{B}$.
    \[
    \|A_n x - y_k^n\| \leqslant \|A x - A_n x\| + \|A_n x - y_k^n\| \leqslant 2 \varepsilon
    \]
\end{itemize}
\end{proof}

% Вам нужен подарок. Приятно дарить бесполезные вещи

\begin{theorem}[Шаудера]
$E$~--- банахово. $A$~--- компактный оператор тогда и только тогда, когда $A^*$~--- компактный оператор.
\end{theorem}

\begin{proof}
Я не знаю, как это доказать. (c) Коновалов
Докажите для случая $E$~--- банахово и линейных ограниченных операторов.
\end{proof}